\section{Introduction}
\label{sec:intro}

Sentiment classification of user reviews is one of the most common text
classification tasks in Korean industry, yet many teams start with only a few
hundred labelled reviews. Two levers are then available: \emph{label more data} or
\emph{use a more expressive model}. Prior work shows that the preferred text
classifier changes with data scale and efficiency requirements: simple linear
models over sparse $n$-gram features are strong at small scale and extremely
cheap~\cite{joulin2017bag}, averaging networks are competitive with syntactic
models~\cite{iyyer2015deep}, and deep models pull ahead only when training sets
are large~\cite{zhang2015character}. Learning curves can be used to predict how
performance grows with data~\cite{kolachina2012prediction}.

This project turns the model choice into a practical data-collection and
deployment decision for Korean reviews.

\paragraph{Research question.}
\emph{For each Korean review dataset, what is the minimum training size among
500, 1{,}000, 2{,}000, 4{,}000 and 8{,}000 reviews at which the best neural sequence
model exceeds the best TF-IDF baseline by at least \crossoverMargin{} macro-F1
point? At that threshold, does the neural model remain competitive in training
time, inference latency and memory use?}

\paragraph{Contributions.}
\begin{itemize}
  \item A controlled learning-curve comparison of four TF-IDF and five neural
        pipelines on three Korean review domains with identical nested training
        subsets, fixed validation/test sets and equal tuning budgets.
  \item A crossover threshold per dataset, with paired per-seed gaps and a
        sensitivity check on how the family winners are selected.
  \item A tokenizer ablation (eojeol, Kiwi morphemes, characters) testing whether the
        conclusion depends on the representation.
  \item A cost profile (fit time, batch-1 latency on CPU and accelerator, peak memory,
        artefact size, model size) and a Pareto analysis at the crossover.
  \item Fully reproducible code: every number in this report is generated from
        the result files (\cref{sec:setup}).
\end{itemize}
